\SourceLeaf{413}%
\par
Ich wähle zur Abscissen-Linie die Gerade \(\mathrm{(II)}\), welche \((4),(5)\) in ihrer Mitte \((6)\) unter rechten Winkeln schneidet, und zähle die Abscissen von \((6)\) an positiv auf der Seite von \((4),(5)\), wo der Winkel von \((4)\) nach \((5)\) unter \(180^0\) erscheint, d.\,i. auf der rechten, wenn man die Winkel von der Linken nach der Rechten wachsen lässt; die Ordinaten mögen in dem Sinne von \((6)\) nach \((5)\) positiv gezählt werden. Auf \(\mathrm{(II)}\) bezeichne ich die Punkte, deren Abscissen
\[c.\operatorname{cotang}\vartheta=n-a,\quad c.\operatorname{cotang}\vartheta'=n,\quad c.\operatorname{cotang}\vartheta''=n+b\]
sind, mit \((1^*),(2^*),(3^*)\); sie sind die Mittelpunkte der drei Kreise, welche beziehungsweise durch \((1),(2),(3)\) und zugleich alle durch \((4)\) und \((5)\) gehen. Die Halbmesser dieser Kreise sind
\[\frac c{\sin\vartheta}=\BareRoot{(cc+(n-a)^2)},\quad\frac c{\sin\vartheta'}=\BareRoot{(cc+nn)},\quad\frac c{\sin\vartheta''}=\BareRoot{(cc+(n+b)^2)}\]
oder wenn man \(\frac c{\sin\vartheta'}=r\) setzt, so werden die beiden andern \(\BareRoot{(rr-2an+aa)},\BareRoot{(rr+2bn+bb)}\). Endlich sei \((7)\) der Durchschnittspunkt von \(\mathrm{(I)}\) und \(\mathrm{(II)}\), \(T\) und \(t\) die Abstände der Punkte \((2)\) und \((2^*)\) von \((7)\), \(\varphi\) der Winkel zwischen gleichnamigen Armen jener Linien, und zwar von \(\mathrm{(I)}\) nach \(\mathrm{(II)}\) in dem gewählten Sinn positiv gemessen. Es ist also die Abscisse von \((7)=n-t\), und folglich sind die Coordinaten der drei Beobachtungsplätze
\[\begin{array}{rcll}(1)&.\ .\ .\ .\ .&n-t+(T-A).\cos\varphi,&(T-A).\sin\varphi\\
(2)&.\ .\ .\ .\ .&n-t+T.\cos\varphi,&T.\sin\varphi\\
(3)&.\ .\ .\ .\ .&n-t+(T+B)\cos\varphi,&(T+B).\sin\varphi\end{array}\]
\par
Die drei unbekannten Grössen \(t,T,\varphi\) werden aber aus folgenden Gleichungen zu bestimmen sein, wenn zur Abkürzung \(x\) für \(\cos\varphi\) geschrieben ist:
\[\begin{aligned}tt+TT-2tTx&=rr\quad.\ .\ .\ .\ .\ .\ [1]\\
(t-a)^2+(T-A)^2-2(t-a)(T-A)x&=rr-2an+aa\\
(t+b)^2+(T+B)^2-2(t+b)(T+B)x&=rr+2bn+bb\end{aligned}\]
\par
Anstatt der beiden letztern gebrauche ich die folgenden, die aus ihrer Subtraction von der ersten hervorgehen
\[\begin{aligned}2at+2AT-2an-AA&=(2At+2aT-2aA)x\quad.\ .\ .\ [2]\\
2bt+2BT-2bn+BB&=(2Bt+2bT+2bB)x\quad.\ .\ .\ [3]\end{aligned}\]
